\documentclass[aspectratio=169,11pt]{beamer}
% Shared Beamer style for practice contest solution walkthroughs.
\usepackage[utf8]{inputenc}
\usepackage[T1]{fontenc}
\usepackage{lmodern}
\usepackage{amsmath,amssymb}
\usepackage{xcolor}
\usepackage{hyperref}
\usepackage{listings}

\definecolor{CPNavy}{HTML}{172033}
\definecolor{CPBlue}{HTML}{2563EB}
\definecolor{CPGreen}{HTML}{16A34A}
\definecolor{CPOrange}{HTML}{F97316}
\definecolor{CPGray}{HTML}{64748B}
\definecolor{CPLight}{HTML}{F8FAFC}
\definecolor{CPCodeBg}{HTML}{F1F5F9}
\definecolor{CPCodeRule}{HTML}{CBD5E1}

\usetheme{default}
\usefonttheme{professionalfonts}
\setbeamertemplate{navigation symbols}{}
\setbeamersize{text margin left=0.65cm,text margin right=0.65cm}
\setbeamertemplate{itemize item}{\color{CPBlue}$\blacktriangleright$}
\setbeamertemplate{itemize subitem}{\color{CPGray}$\triangleright$}

\setbeamercolor{normal text}{fg=CPNavy,bg=white}
\setbeamercolor{frametitle}{fg=white,bg=CPNavy}
\setbeamercolor{title}{fg=white,bg=CPNavy}
\setbeamercolor{subtitle}{fg=CPLight,bg=CPNavy}
\hypersetup{colorlinks=true,urlcolor=CPBlue}

\setbeamertemplate{frametitle}{%
  \nointerlineskip%
  \begin{beamercolorbox}[wd=\paperwidth,ht=0.88cm,dp=0.22cm,leftskip=0.55cm,rightskip=0.35cm]{frametitle}
    \usebeamerfont{frametitle}\insertframetitle\hfill\small\insertframenumber/\inserttotalframenumber
  \end{beamercolorbox}%
}

\setbeamertemplate{footline}{%
  \leavevmode%
  \hbox{%
    \begin{beamercolorbox}[wd=.58\paperwidth,ht=2.4ex,dp=1ex,leftskip=.5cm]{author in head/foot}%
      \color{CPGray}\insertshorttitle
    \end{beamercolorbox}%
    \begin{beamercolorbox}[wd=.42\paperwidth,ht=2.4ex,dp=1ex,rightskip=.5cm plus1fil]{date in head/foot}%
      \hfill\color{CPGray}\insertshortauthor
    \end{beamercolorbox}%
  }%
}

\setbeamertemplate{title page}{%
  \vspace*{1.25cm}
  \begin{beamercolorbox}[wd=\paperwidth,sep=0.65cm,leftskip=0.15cm,rightskip=0.15cm]{title}
    {\color{white}\Huge\bfseries\inserttitle\par}
    \vspace{0.35cm}
    {\color{CPLight}\Large\insertsubtitle\par}
    \vspace{0.6cm}
    {\color{CPLight}\large\insertauthor\par}
  \end{beamercolorbox}
  \vfill
}

\newcommand{\sectiontag}[1]{\textbf{\color{CPBlue}#1}}
\newenvironment{tightitemize}{\begin{itemize}\small\setlength{\itemsep}{0.18cm}}{\end{itemize}}
\lstdefinestyle{cpcode}{
  language=Python,
  basicstyle=\ttfamily\tiny,
  keywordstyle=\color{CPBlue}\bfseries,
  commentstyle=\color{CPGray}\itshape,
  stringstyle=\color{CPGreen},
  backgroundcolor=\color{CPCodeBg},
  frame=single,
  framerule=0.25pt,
  rulecolor=\color{CPCodeRule},
  showstringspaces=false,
  columns=fullflexible,
  keepspaces=true,
  breaklines=true,
  breakatwhitespace=false,
  tabsize=4,
  xleftmargin=0.08cm,
  xrightmargin=0.08cm,
  aboveskip=0.08cm,
  belowskip=0.08cm
}


\title[ceiling]{Ceiling Function}
\subtitle{Day 2 solution walkthrough}
\author[Solution Deck]{Competitive Programming Bootcamp}
\date{}

\begin{document}

\begin{frame}[plain]
  \titlepage
\end{frame}

% BEGIN GENERATED STATEMENT INTERPRETATION
\begin{frame}{Interpret the Word Problem}
\small
\sectiontag{Statement excerpts:}
\begin{tightitemize}
  \item \textit{``store them in a binary search tree''}
  \item \textit{``trees of the same shape''}
\end{tightitemize}

\vspace{0.18cm}
\sectiontag{Programming interpretation:}
\begin{tightitemize}
  \item Each prototype is a sequence of insertions into a BST.
  \item The actual numeric labels do not matter after insertion; only left/right shape matters.
  \item Serialize each shape into a tuple and count unique tuple shapes.
\end{tightitemize}
\end{frame}
% END GENERATED STATEMENT INTERPRETATION







\begin{frame}{Problem Model}
\small
\sectiontag{Textbook connection:} CP4 2.3.3 Balanced BST / 3.2 Search Modeling

\vspace{0.25cm}
\sectiontag{Core technique:} BST construction and structural hashing

\vspace{0.25cm}
\begin{tightitemize}
  \item Each sequence is inserted into a binary search tree in input order.
  \item Only the shape of the resulting tree matters.
  \item The answer is the number of distinct tree shapes.
\end{tightitemize}
\end{frame}

\begin{frame}{How to Recognize the Approach}
\begin{tightitemize}
  \item Different values can create the same comparison structure.
  \item If two trees have the same left/right recursive shape, they are equivalent.
  \item A set can store canonical shape descriptions and remove duplicates.
\end{tightitemize}
\end{frame}

\begin{frame}{Algorithm}
\begin{tightitemize}
  \item For each sequence, create a root from the first value.
  \item Insert remaining values using ordinary BST rules.
  \item Convert the finished tree to a recursive tuple shape.
  \item Use None for empty children and ignore stored values.
  \item Insert the tuple into a set and print the set size.
\end{tightitemize}
\end{frame}

\begin{frame}{Walkthrough of the Key State}
\begin{tightitemize}
  \item get\_shape(root) returns (left\_shape, right\_shape).
  \item Leaf nodes become (None, None).
  \item Two equal tuples mean the same exact branching structure.
\end{tightitemize}
\end{frame}

\begin{frame}{Why the Algorithm Is Correct}
\begin{tightitemize}
  \item BST insertion is simulated exactly for every sequence.
  \item The tuple representation records every left and right child position.
  \item Ignoring values is valid because the problem compares only tree shape.
\end{tightitemize}
\end{frame}

\begin{frame}{Complexity and Implementation Notes}
\small
\sectiontag{Complexity:} O(N * K\textasciicircum{}2) worst-case time for N sequences of length K; O(N * K) memory for stored shapes.

\vspace{0.3cm}
\begin{tightitemize}
  \item The tree may be unbalanced, but the input bounds make this acceptable.
  \item Tuples are hashable, so they can be stored directly in a Python set.
\end{tightitemize}
\end{frame}



% BEGIN GENERATED CODE WALKTHROUGH
\begin{frame}[fragile]{Code: Represent a BST Node}
\scriptsize
\sectiontag{Source lines:} \texttt{ceiling.py:32--40}

\vspace{0.08cm}
\begin{columns}[T,totalwidth=\textwidth]
\begin{column}{0.58\textwidth}
\begin{lstlisting}[style=cpcode]
import sys

class Node:
    def __init__(self, value):
        self.value = value
        self.left = None
        self.right = None
\end{lstlisting}
\end{column}
\begin{column}{0.38\textwidth}
\sectiontag{What this chunk does}
\begin{tightitemize}
  \item Each node stores the inserted value and pointers to left and right children.
  \item None children are important because they distinguish tree shapes.
\end{tightitemize}
\end{column}
\end{columns}
\end{frame}

\begin{frame}[fragile]{Code: Insert by BST Rules}
\scriptsize
\sectiontag{Source lines:} \texttt{ceiling.py:42--55}

\vspace{0.08cm}
\begin{columns}[T,totalwidth=\textwidth]
\begin{column}{0.58\textwidth}
\begin{lstlisting}[style=cpcode]
def insert(root, value):
    if value < root.value:
        if root.left is None:
            root.left = Node(value)
        else:
            insert(root.left, value)

    else:
        if root.right is None:
            root.right = Node(value)
        else:
            insert(root.right, value)
\end{lstlisting}
\end{column}
\begin{column}{0.38\textwidth}
\sectiontag{What this chunk does}
\begin{tightitemize}
  \item Smaller values go left; all other values go right.
  \item Insertion order is what creates the final shape.
  \item The code recurses until it finds the empty child position.
\end{tightitemize}
\end{column}
\end{columns}
\end{frame}

\begin{frame}[fragile]{Code: Serialize the Shape}
\scriptsize
\sectiontag{Source lines:} \texttt{ceiling.py:57--64}

\vspace{0.08cm}
\begin{columns}[T,totalwidth=\textwidth]
\begin{column}{0.58\textwidth}
\begin{lstlisting}[style=cpcode]
def get_shape(root):
    if root is None:
        return None

    return (get_shape(root.left), get_shape(root.right))
\end{lstlisting}
\end{column}
\begin{column}{0.38\textwidth}
\sectiontag{What this chunk does}
\begin{tightitemize}
  \item Empty children become None.
  \item A real node becomes a pair containing the shapes of its left and right subtrees.
  \item Values are intentionally ignored, so equal structures produce equal tuples.
\end{tightitemize}
\end{column}
\end{columns}
\end{frame}

\begin{frame}[fragile]{Code: Count Unique Shapes}
\scriptsize
\sectiontag{Source lines:} \texttt{ceiling.py:66--89}

\vspace{0.08cm}
\begin{columns}[T,totalwidth=\textwidth]
\begin{column}{0.58\textwidth}
\begin{lstlisting}[style=cpcode]
def main():
    input = sys.stdin.buffer.readline

    n, k = map(int, input().split())

    shapes = set()

    for _ in range(n):
        values = list(map(int, input().split()))

        root = Node(values[0])

        for value in values[1:]:
            insert(root, value)

        shapes.add(get_shape(root))

    print(len(shapes))

main()
\end{lstlisting}
\end{column}
\begin{column}{0.38\textwidth}
\sectiontag{What this chunk does}
\begin{tightitemize}
  \item Each prototype builds its own BST from the given sequence.
  \item The serialized shape is inserted into a set, which removes duplicates.
  \item The number of set entries is the answer.
\end{tightitemize}
\end{column}
\end{columns}
\end{frame}
% END GENERATED CODE WALKTHROUGH

\begin{frame}{Source}
\small
\sectiontag{Solution file:} \texttt{practice\_contests/day\_2/ceiling.py}

\vspace{0.3cm}
\sectiontag{Problem:} \url{https://open.kattis.com/problems/ceiling}

\vspace{0.45cm}
Use this deck with the source code open beside it. The slides explain the reasoning; the code shows the exact input handling and output formatting.
\end{frame}

\end{document}
